\begin{tikzpicture}[>=Latex, line cap=round, line join=round, font=\small,
  x={(-1cm,-0.65cm)}, y={(1cm,0cm)}, z={(0cm,1cm)}]
\coordinate (origen) at (0,0,0);
\coordinate (verticex) at (2.7,0,0);
\coordinate (verticey) at (0,3.6,0);
\coordinate (verticez) at (0,0,3.4);
\filldraw[fill=gray!8, draw=gray, dashed] (origen) -- (verticey) -- (verticez) -- cycle;
\filldraw[fill=blue!12, draw=blue!55!black, thick] (origen) -- (verticex) -- (verticey) -- cycle;
\filldraw[fill=orange!14, draw=orange!70!black, thick] (origen) -- (verticex) -- (verticez) -- cycle;
\filldraw[fill=green!12, fill opacity=0.65, draw=green!40!black, thick] (verticex) -- (verticey) -- (verticez) -- cycle;
\draw[dashed, gray] (origen) -- (verticey) (origen) -- (verticez);
\draw[->, gray] (verticex) -- (3.4,0,0) node[below] {$x$};
\draw[->, gray] (verticey) -- (0,4.4,0) node[right] {$y$};
\draw[->, gray] (verticez) -- (0,0,4.2) node[above] {$z$};
\node[below right] at (origen) {$O$};
\node[below left] at (verticex) {$A$};
\node[below] at (verticey) {$B$};
\node[left] at (verticez) {$C$};
\coordinate (caraxy) at (0.9,1.2,0);
\coordinate (carazx) at (0.9,0,1.13);
\coordinate (carainclinada) at (0.9,1.2,1.13);
\draw[->, very thick, blue!70!black] (caraxy) -- (0.9,1.2,-1.45) node[below] {$d\vec F_z$};
\draw[->, very thick, orange!85!black] (carazx) -- (0.9,-1.7,1.13) node[left] {$d\vec F_y$};
\draw[->, very thick, red!70!black] (carainclinada) -- (0.9,2.9,2.9) node[right] {$d\vec F$};
\node[green!30!black] at (0.4,2.45,0.6) {$dS$};
\fill (caraxy) circle (1.7pt);
\fill (carazx) circle (1.7pt);
\fill (carainclinada) circle (1.7pt);
\node[align=left, anchor=west, font=\footnotesize] at (0,4.6,2) {$OAB$: plano $xy$, área $dS_z$\\$OAC$: plano $zx$, área $dS_y$\\$OBC$: plano $yz$, área $dS_x$\\$ABC$: cara inclinada, área $dS$};
\node[align=center, font=\footnotesize] at (0,0.7,-3) {$d\vec F_x$ actúa sobre la cara oculta $OBC$.\\Las fuerzas se dibujan esquemáticamente; no tienen que ser normales a las caras.};
\end{tikzpicture}
